\chapter{Implementation Reference}\label{app:apiref}

This appendix documents the implementation-level interface details referenced from
\cref{ch:impl}: the directory server's REST API, the privileged testing interface, the
Signal Desktop integration hooks and user-interface events, the repository layout, and the
runner commands. None of it is needed to follow the design (\cref{ch:design}) or the
evaluation (\cref{ch:eval}); it is provided for reproducibility and for readers who wish to
run the system.

% ----------------------------------------------------------------------------
\section{Directory Server REST API}\label{sec:apiref-rest}

The standalone server exposes a public API used by clients and auditors. Requests and
responses are JSON, and byte strings are hex-encoded. The verified-lookup endpoints return the
value together with the prefix proof, the log-inclusion proof, and the signed tree head, so
that the client can run the full verification chain (\cref{ssec:design-verify}).
\Cref{tab:apiref-public} lists these public endpoints.

\begin{table}[ht]
\centering
\caption{Public directory-server endpoints.}
\label{tab:apiref-public}
\setlength{\extrarowheight}{2pt}
\footnotesize
\begin{tabular}{@{}lp{3.9cm}p{7.6cm}@{}}
\toprule
Method & Path & Purpose \\
\midrule
GET  & \texttt{/sth} & current signed tree head \\
GET  & \texttt{/search} & verified lookup by label; returns value, prefix proof, log-inclusion proof, and the signed STH \\
GET  & \texttt{/search-by-user-id} & lookup by user identifier, with the VRF proof attached \\
GET  & \texttt{/consistency} & consistency proof between two log sizes (\texttt{from}, \texttt{to}) \\
GET  & \texttt{/epoch} & current epoch number \\
GET  & \texttt{/log-root-at} & log root at a given tree size \\
GET  & \texttt{/epoch-delta} & a committed epoch's declared update set and committed directory root, so that a networked auditor can replay the transition (\cref{alg:auditor}) \\
GET  & \texttt{/vrf-public-key} & operator VRF public key \\
GET  & \texttt{/sth-public-key} & operator STH signing public key \\
POST & \texttt{/register} & register a user's key (server derives the VRF label) \\
POST & \texttt{/update-raw} & write a client-derived (dual-index) label$\to$value; the public surface is write-once, so an occupied label is rejected with \texttt{409 Conflict} (\cref{ssec:design-detection}) \\
POST & \texttt{/commit} & close the current epoch and sign a new STH \\
GET  & \texttt{/health} & liveness probe \\
GET  & \texttt{/dashboard}, \texttt{/internals}, \texttt{/integration}, \texttt{/threats} & served visualization and threat-model pages \\
\bottomrule
\end{tabular}
\end{table}

% ----------------------------------------------------------------------------
\FloatBarrier
\section{Privileged Testing Interface}\label{sec:apiref-admin}

A separate administrative API, bound to an internal-only port, is used exclusively by the
adversarial scenario suite to stage attacks; it is not part of the protocol. In the
prototype these endpoints are unauthenticated, and a production deployment would place them
behind account-layer authorization. \Cref{tab:apiref-admin} lists them.

\begin{table}[ht]
\centering
\caption{Privileged testing endpoints for internal use.}
\label{tab:apiref-admin}
\setlength{\extrarowheight}{2pt}
\small
\begin{tabular}{@{}llp{7.8cm}@{}}
\toprule
Method & Path & Purpose \\
\midrule
POST & \texttt{/admin/update} & overwrite a slot's value, bypassing the public write-once policy; the overwrite is still declared in the epoch delta as a rotation (stages A1/A5 substitution) \\
POST & \texttt{/admin/delete} & remove a committed leaf \emph{without} declaring it in the epoch delta, so the auditor's transition replay fails (stages the suppression-by-deletion scenario) \\
POST & \texttt{/admin/commit} & commit an epoch \\
POST & \texttt{/admin/equivocate} & sign a second, divergent head at one epoch (stages A4) \\
POST & \texttt{/admin/stop-equivocate} & clear the forked head \\
POST & \texttt{/admin/reset} & rebuild the server from its fixed seed with a fresh genesis epoch, for repeatable demonstrations \\
\bottomrule
\end{tabular}
\end{table}

% ----------------------------------------------------------------------------
\FloatBarrier
\section{Signal Desktop Integration Hooks and Events}\label{sec:apiref-hooks}

\Cref{tab:apiref-hooks} maps the integration concepts of \cref{sec:impl-signal} to the
corresponding prototype functions, and \cref{tab:apiref-events} maps the
user-interface verdicts to their DOM event names. The live client attaches only at key upload and
peer-key retrieval and records peer identity through a periodic sweep. The session-establishment and
pre-key-bundle rows are documentation-only integration sketches that mark where a production binding
of the live X3DH secret would attach.

\begin{table}[ht]
\centering
\caption{Integration hooks.}
\label{tab:apiref-hooks}
\small
\begin{tabular}{@{}lll@{}}
\toprule
Lifecycle point & Prototype function & Layer \\
\midrule
key upload & \texttt{registerKeys} & TypeScript \\
peer-key retrieval & \texttt{getKeysForServiceId} & TypeScript \\
session establishment & \texttt{sessionEstablished} & TypeScript \\
pre-key-bundle processing & \texttt{hook\_process\_prekey\_bundle} & Rust (libsignal) \\
\bottomrule
\end{tabular}
\end{table}

\begin{table}[ht]
\centering
\caption{User-interface verdict events.}
\label{tab:apiref-events}
\small
\begin{tabular}{@{}ll@{}}
\toprule
Meaning & DOM event \\
\midrule
warning / man-in-the-middle detected (red) & \texttt{kt:warning} \\
key verified (green) & \texttt{kt:verified} \\
server unreachable (amber) & \texttt{kt:unreachable} \\
demonstration: inject a ratchet fork & \texttt{kt:cea-inject} \\
\bottomrule
\end{tabular}
\end{table}

% ----------------------------------------------------------------------------
\FloatBarrier
\section{Repository Layout and Runner Commands}\label{sec:apiref-repo}

The implementation is organized as follows:

\begin{itemize}
  \item \texttt{crate/}: the Rust \texttt{duet} crate, with library modules
  (\cref{tab:impl-modules}), the standalone server and auditor binaries, and the benchmark
  harness.
  \item \texttt{demo/}: the server, auditor, and attacker orchestration together with
  the served visualization and threat-model pages.
  \item \texttt{signal-desktop/ts/kt/}: the embedded TypeScript client, whose shared
  verifier modules are mirrored byte-for-byte by the standalone \texttt{signal-desktop-kt/}
  package, while integration wiring and test scaffolding differ between the two.
\end{itemize}

The system is driven by single commands: \texttt{cargo test} and \texttt{cargo bench}, run from
\texttt{crate/}, for the Rust tests and benchmarks, and the runner script
\texttt{demo/run.ps1} for the multi-actor demonstration. That script's native subcommands are
\texttt{server}, \texttt{auditor}, \texttt{build}, \texttt{client1}, \texttt{client2},
\texttt{attack}, \texttt{fork}, \texttt{delete}, \texttt{simulate-delete}, \texttt{stop-fork},
\texttt{stop}, and \texttt{reset}; \texttt{client1} and \texttt{client2} launch the two Signal
instances the two-client demonstration needs. A containerized path
(\texttt{up}, \texttt{logs}, \texttt{down}) runs the server, auditor, and attacker under Docker
instead. The networked-in-app mode is enabled by setting a server URL in
the client's environment before launch.

% ----------------------------------------------------------------------------
\FloatBarrier
\section{Integration and Flow Diagrams}\label{sec:apiref-diagrams}

These diagrams accompany the client integration of \cref{sec:impl-signal}. \Cref{fig:impl-signal}
maps the runtime arrangement of the embedded client and the three processes it drives,
\cref{fig:impl-verify} traces the matching branch of the networked
\texttt{RemoteKTClient.detectMitm} verification path, and \cref{fig:impl-e2e} traces the normal and
A1b attack workflows across the four roles.

\begin{figure}[ht]
\centering
% fig-impl-signal.tex — extracted from ch5-implementation.tex
% Label: \label{fig:impl-signal}
\begin{tikzpicture}[>=Stealth,font=\footnotesize,
   box/.style={draw,rounded corners,align=center,minimum height=8mm},
   sub/.style={draw,rounded corners,align=center,minimum height=6mm,font=\scriptsize,fill=white}]
  \node[box,minimum width=6.8cm,minimum height=3.8cm,fill=black!4] (app) at (0,0) {};
  \node[font=\small] at (0,1.6) {\textbf{Signal Desktop (Electron)}};
  \node[sub,minimum width=6.0cm] (hooks) at (0,0.85)
        {hooks: key upload \,|\, key lookup\\session establishment \,|\, safety-number change};
  \node[sub,minimum width=2.2cm] (client) at (-1.75,-0.75) {KT client\\(verifier)};
  \node[sub,minimum width=2.0cm] (ui) at (1.8,-0.75) {Security UI\\(badge / banner)};
  \draw[->] (client) -- node[above,font=\scriptsize]{verdicts} (ui);
  \node[box,minimum width=2.4cm,fill=black!8] (srv) at (7.0,0.4) {Rust server};
  \node[box,minimum width=2.0cm] (aud) at (5.8,-2.0) {auditor};
  \node[box,minimum width=2.0cm,draw=red!60!black] (atk) at (8.2,-2.0) {attacker};
  \draw[<->] (app.east) -- node[above,font=\scriptsize]{HTTP} (srv.west);
  \draw[->,dashed] (aud) -- node[left,font=\scriptsize]{poll STH} (srv);
  \draw[->,dashed,red!60!black] (atk) -- node[right,font=\scriptsize]{stage attack} (srv);
\end{tikzpicture}

\caption[Signal Desktop integration points.]{Signal Desktop integration. The embedded TypeScript
KT client communicates over HTTP with the Rust directory server. It verifies keys at the key-upload
and peer-key-retrieval hooks and reports verdicts to the badge and banner. An independent auditor polls the server. The
scripted attacker uses a privileged testing interface. \Cref{tab:apiref-hooks,tab:apiref-events} list the hook functions
and verdict-event names.}
\label{fig:impl-signal}
\end{figure}

\begin{figure}[ht]
\centering
% fig-impl-verify.tex — extracted from ch5-implementation.tex
% Label: \label{fig:impl-verify}
\begin{tikzpicture}[>=Stealth,font=\footnotesize,
   hd/.style={draw,rounded corners,minimum width=3.4cm,minimum height=7mm,align=center,fill=black!6},
   act/.style={draw,rounded corners,fill=black!4,align=center,font=\scriptsize,inner sep=3pt}]
  % lifeline heads
  \node[hd] (c) at (0,0) {\textbf{Client}\\{\scriptsize(\texttt{RemoteKTClient})}};
  \node[hd] (s) at (7.4,0) {\textbf{Server}};
  \draw[dashed] (c.south) -- (0,-6.6);
  \draw[dashed] (s.south) -- (7.4,-6.6);
  % read own slot
  \draw[->] (0,-0.9) -- node[above,font=\scriptsize]{read own slot at $\ell_{\mathsf{self}}$} (7.4,-0.9);
  % response (label above the arrow, shortened, so it clears the lifeline + block)
  \draw[<-] (0,-2.0) -- node[above,font=\scriptsize]{value $+$ proof chain $+$ signed STH} (7.4,-2.0);
  % client-side verification
  \node[act] at (0,-3.05)
     {verify STH signature\\recompute + verify prefix proof\\verify log inclusion + label binding\\anchor head (append-only)};
  % read warning slot (only if own slot matched)
  \draw[->] (0,-4.25) -- node[above,font=\scriptsize]{if own slot matches: read $\ell_{\mathsf{warn}}$} (7.4,-4.25);
  \draw[<-] (0,-5.1) -- node[below,font=\scriptsize]{proven (non-)inclusion} (7.4,-5.1);
  % verdict
  \node[act,fill=black!8] at (0,-6.15) {return:\ \texttt{null} (safe) / \texttt{MitmWarning}};
\end{tikzpicture}

\caption[Owner-slot and warning-slot verification.]{Matching branch of
\texttt{RemoteKTClient.detectMitm}. The client verifies its own dual-index slot against a signed
head anchored to the append-only log. Only a matching owner-slot value proceeds to the proof-verified
warning-slot lookup. The live monitor handles registration and warning publication separately.}
\label{fig:impl-verify}
\end{figure}

\begin{figure}[p]
\centering
% fig-impl-e2e.tex — end-to-end sequence across the four roles: the normal path
% (top) and one compact attack branch (bottom). Label: \label{fig:impl-e2e}
% Two client lifelines (Alice, Bob), each with its verifier embedded; the KT
% server; and the auditor. Every proof-verified lookup is drawn separately.
\begin{tikzpicture}[>=Stealth,font=\scriptsize,
  hd/.style={draw,rounded corners,fill=black!6,align=center,minimum height=9mm,minimum width=21mm,inner sep=3pt},
  act/.style={draw,fill=black!4,align=center,font=\tiny,inner sep=2.5pt,rounded corners,text width=2.9cm},
  ract/.style={act,draw=red!60!black},
  gact/.style={act,draw=green!45!black,fill=green!7},
  m/.style={->,shorten >=1pt},
  red/.style={->,red!65!black,shorten >=1pt},
  grn/.style={->,green!50!black,shorten >=1pt},
  mlbl/.style={fill=white,inner sep=1.5pt,font=\scriptsize}]

  % lifelines (x positions): Alice 0, Bob 3.4, KT server 7.0, Auditor 10.6
  \foreach \x/\name/\id in {0/{Alice\\{\tiny(client + verifier)}}/A, 3.4/{Bob\\{\tiny(client + verifier)}}/B, 7.0/{KT server}/S, 10.6/Auditor/U}{
    \node[hd] (\id) at (\x,0) {\name};
    \draw[dashed,black!45] (\x,-0.55) -- (\x,-11.4);
  }
  \node[font=\itshape,text=black!55,anchor=west,fill=white,inner sep=1.5pt] at (-0.95,-0.9) {normal operation};

  % ---- normal path ----
  \draw[m] (0,-1.5) -- node[above,mlbl]{register Alice's key at $\ell_A$; commit epoch} (7.0,-1.5);
  \draw[m] (3.4,-2.3) -- node[above,mlbl]{search peer slot $\ell_A$} (7.0,-2.3);
  \draw[m] (7.0,-3.0) -- node[above,mlbl]{value $+$ proof chain $+$ STH} (3.4,-3.0);
  \node[act,anchor=north] at (3.4,-3.35) {three label-bound verified lookups:\\peer slot $\ell_A$, own slot $\ell_B$,\\warning slot $\ell_{\mathsf{warn}}$ (empty)};
  \node[gact,anchor=north] at (3.4,-5.15) {\textsc{safe} (green)};
  \draw[m] (10.6,-4.1) -- node[above,mlbl]{poll STH $+$ delta} (7.0,-4.1);
  \node[act,anchor=north] at (10.6,-4.45) {replay transition:\\consistent};

  % ---- divider ---- (lifelines span x=0..10.6; midpoint 5.3)
  \draw[dashed,red!55!black] (-0.95,-6.1) -- (11.55,-6.1);
  \node[font=\itshape,text=red!60!black,fill=white,inner sep=2pt] at (5.3,-6.42)
    {under attack (A1b: overwrite of the immutable dual-index slot $\ell_A$)};

  % ---- attack branch ----
  % server overwrites the value (box centred on the server lifeline)
  \node[ract,anchor=north,text width=2.5cm] at (7.0,-6.85) {server overwrites the value at $\ell_A$ with Mallory's key};
  % auditor detects the illegal rotation (kept clear of the server box, right of it)
  \draw[red] (10.6,-8.05) -- node[above,mlbl,font=\tiny,pos=0.4]{poll STH $+$ delta} (7.0,-8.05);
  \node[ract,anchor=north] at (10.6,-8.4) {replay: declared immutable rotation $\to$ reject};
  % Alice's monitor detects the overwrite of her own slot
  \draw[m] (0,-8.7) -- node[above,mlbl]{monitor tick: search own slot $\ell_A$} (7.0,-8.7);
  \draw[red] (7.0,-9.3) -- node[above,mlbl]{Mallory's key} (0,-9.3);
  \node[ract,anchor=north] at (0,-9.65) {own-slot mismatch:\\pinned key $\ne$ Mallory};
  % Alice writes the warning
  \draw[red] (0,-10.55) -- node[above,mlbl]{write warning slot $\ell_{\mathsf{warn}}$} (7.0,-10.55);
  % Bob's monitor reads the warning
  \draw[red] (7.0,-11.05) -- node[above,mlbl]{Bob monitor: warning present} (3.4,-11.05);
  \node[ract,anchor=north] at (3.4,-11.4) {\textsc{warning} (red)};
\end{tikzpicture}

\caption[Normal and A1b attack paths across the four roles.]{Normal and A1b attack paths across
the two clients, directory, and auditor. In the lower panel, Alice detects an own-slot mismatch and
writes a warning that Bob later reads. The auditor separately rejects the immutable-slot rotation.}
\label{fig:impl-e2e}
\end{figure}

In \cref{fig:impl-e2e}, Alice and Bob each run an embedded verifier, while the KT server and auditor
execute as separate processes. In the normal path Alice registers her key and commits. Bob's verifier
then performs three label-bound, proof-verified lookups (the peer slot $\ell_A$ it retrieves Alice
from, its own slot $\ell_B$, and the shared warning slot $\ell_{\mathsf{warn}}$, which is verifiably
empty) and returns \texttt{SAFE}, while the auditor polls and replays to \textsc{Consistent}. In the
A1b branch the server overwrites Alice's own dual-index slot $\ell_A$ (which is also Bob's peer slot).
Alice's monitor reads the substituted value against her pinned key, finds a mismatch, and writes the
shared warning slot. Bob's monitor then reads that warning and both clients go red. Independently, the
auditor's transition replay rejects the declared immutable rotation, so detection does not rest on the
client path alone. The live run in \cref{ssec:impl-demo} shows this inside the real interface.

% ----------------------------------------------------------------------------
\FloatBarrier
\section{Demonstration Gallery}\label{sec:apiref-demo}

The supporting captures for the demonstration of \cref{ssec:impl-demo} show the
honest baseline before any attack, the attacker-side terminals that stage each scenario, the
global-VRF and continuous-mode detection fronts, and the auditor outputs. The composite two-client
warning capture that carries the RQ4 evidence remains in \cref{ssec:impl-demo}
(\cref{fig:impl-screens}). In that capture the two DevTools consoles record the bidirectional
detection at one shared slot. The affected client logs a warning-slot write, and the peer logs a read
at the identical label. The honest baseline is shown in \cref{fig:demo-baseline}, the attacker
terminals in \cref{fig:demo-cmd}, the staged equivocation in \cref{fig:demo-equivocation}, the
global-VRF and CEA fronts in \cref{fig:impl-screens-fronts}, and the auditor outputs in
\cref{fig:impl-auditor}.

\begin{figure}[ht]
\centering
\screenshotplaceholder{0.49\textwidth}{figures/ui-baseline-self.png}\hfill
\screenshotplaceholder{0.49\textwidth}{figures/ui-baseline-contact.png}
\caption[Honest baseline (scenario A0).]{Honest baseline for scenario A0. Two linked Signal Desktop
clients verify keys against a clean directory and display no warning banner. This state provides the
reference for the later attack captures.}
\label{fig:demo-baseline}
\end{figure}

\begin{figure}[ht]
\centering
\screenshotplaceholder{0.49\textwidth}{figures/cmd-dualindex.png}\hfill
\screenshotplaceholder{0.49\textwidth}{figures/cmd-vrf.png}
\caption[Attacker terminals staging the two substitutions.]{Attacker terminals staging two
substitutions. The left capture shows \texttt{attack.mjs --label <hex>} overwriting the
per-conversation lookup slot. The right capture shows \texttt{attack.mjs} with a target user ID
overwriting the VRF-derived identity entry.}
\label{fig:demo-cmd}
\end{figure}

\begin{figure}[ht]
\centering
\screenshotplaceholder{0.62\textwidth}{figures/cmd-equivocation.png}
\caption[Staging an equivocation (scenario A4).]{Staging an equivocation in scenario A4. The
privileged testing endpoint signs a second, divergent tree head at the same epoch. The
cross-vantage auditor detects the divergence (\cref{fig:impl-auditor}, right).}
\label{fig:demo-equivocation}
\end{figure}

\begin{figure}[ht]
\centering
\screenshotplaceholder{0.49\textwidth}{figures/ui-vrf.png}\hfill
\screenshotplaceholder{0.49\textwidth}{figures/ui-cea.png}
\caption[Global-VRF and CEA detection paths.]{Global-VRF and CEA detection paths. The left capture
shows the account-wide monitor detecting replacement of the VRF-derived identity entry in scenario
A1a. Its console reports \emph{the server replaced it (MitM, attack A1a)}. The right capture shows
continuous CEA mode detecting a conflicting per-epoch fingerprint after ratchet-key injection.
Both paths raise the red warning banner.}
\label{fig:impl-screens-fronts}
\end{figure}

\begin{figure}[ht]
\centering
\screenshotplaceholder{0.49\textwidth}{figures/auditor-rotations.png}\hfill
\screenshotplaceholder{0.49\textwidth}{figures/auditor-fork.png}
\caption[Auditor outputs for A1a, A1b, and A4.]{Auditor outputs for three attacks. In the left
capture, A1a is a permitted declared rotation of a \textsc{mutable} identity slot, so the auditor
reports \textsc{Consistent} and the client self-check provides detection. The same capture shows
A1b causing an \textsc{ImmutableRotation} transition-audit failure. In the right capture, A4 raises
\textsc{alarm: fork}.}
\label{fig:impl-auditor}
\end{figure}

% ============================================================================
\section{Implementation Detail and Engineering Record}\label{sec:apiref-engineering}

This section preserves the module, server, integration, and engineering detail summarized in
\cref{ch:impl}.

\subsection{Directory State, Cryptography, and Server}\label{ssec:apiref-eng-core}

The directory's state lives entirely in memory, in the cooperating structures that parallel the
authenticated-data-structure stack of the design (\cref{ssec:design-structures}). The base
\texttt{hashing} module fixes the RFC~6962 leaf and node domain tags shared by both structures.
The \texttt{merkle\_log} module maintains the append-only log as a vector of hashes and implements
the RFC~9162 inclusion and consistency verifiers, each proof costing time logarithmic in the
number of entries. The \texttt{prefix\_tree} module implements the directory itself as a
path-compressed sparse Merkle tree, held as a pointer-based structure so that only the nodes
actually populated are materialized. Each compressed leaf lazily memoizes the top of its
empty-sibling hash chain, so that serving a proof whose copath meets a compressed leaf does not
recompute the up-to-256-level chain on every lookup. A single \texttt{Arc<Mutex<KTServer>>} protects the server state, so reads and writes are serialized.
The prototype does not shard or replicate this state.

The \texttt{vrf} module implements the
ECVRF-EDWARDS25519-SHA512-TAI construction of RFC~9381 (suite \texttt{0x03}). The \texttt{dual\_index} module
derives the three labels $\ell_{\mathsf{self}}$, $\ell_{\mathsf{peer}}$, and $\ell_{\mathsf{warn}}$
from the shared secret by HKDF-SHA256 (RFC~5869), exactly as specified in
\cref{ssec:design-labels}, and the \texttt{cea} module implements the continuous example construction
of the ACKA framework~\cite{dowling2023acka}. The crates
\texttt{sha2}, \texttt{hkdf}, \texttt{ed25519-dalek}, and \texttt{curve25519-dalek} supply the standard
primitives, and the ECVRF of RFC~9381 is the one hand-written primitive, checked against the published
suite-\texttt{0x03} test vectors.

A standalone HTTP server, built on the
\texttt{axum} framework over the asynchronous \texttt{tokio} runtime, wraps the \texttt{KTServer}
and serves the lookup, monitoring, and consistency operations as JSON, alongside a privileged
testing interface used only by the adversarial scenario suite and the scripted attacker. An independent
auditor binary, a small \texttt{ureq} HTTP client, verifies signed-head growth and replays epoch
transitions along its own view. The server enforces the write-once policy in its raw-label
update path, so a write to an already-occupied dual-index, CEA, or warning label is rejected (the public
\texttt{/update-raw} endpoint answers \texttt{409 Conflict}), while the privileged interface exposes
an explicit unchecked overwrite so the adversarial suite can play a malicious operator. Every commit
freezes a \emph{declared update set} that the server publishes at an \texttt{/epoch-delta} endpoint,
which the auditor consumes (\cref{fig:impl-transition-audit}) by replaying each epoch's declared set
against its own replica and accepting only if the replayed root equals the committed root. This
corresponds to the third-party auditor role in the IETF KeyTrans protocol~\cite{ietf-keytrans-protocol}.
Independently of any auditor, the \texttt{monitor} module provides a writer-side persistence check that
retains the inclusion proof and signed tree head from the epoch in which a warning commits, so that a
later removal yields a bundle verifying against the server's own signed heads (\cref{alg:detect}).
Authenticated ownership of raw-label writes is not implemented, so without account-layer authorization a
malicious client could occupy an unclaimed peer slot before its legitimate owner writes to it. The one
core cost that is not sublinear is the epoch commit, which snapshots the tree by a deep clone.

\subsection{Signal Desktop Integration Wiring}\label{ssec:apiref-eng-wiring}

The integration uses key-upload and peer-key-retrieval hooks together with a background monitor.
Peer identity material is recorded during key retrieval and through a periodic identity-record
sweep. No live session-establishment hook exposes the X3DH shared secret. The
peer-key-retrieval hook runs the full verification chain
(\cref{ssec:design-verify}) on a peer's key at the retrieval boundary and reports the verdict without
blocking session establishment. Verdicts appear in the user
interface through a conversation-header badge and a
banner, driven by DOM events. The specific hook functions and event names are
listed in \cref{tab:apiref-hooks,tab:apiref-events}.
\Cref{fig:impl-verify} illustrates the matching branch of the networked \texttt{RemoteKTClient.detectMitm}
path: the client reads its own dual-index lookup slot $\ell_{\mathsf{self}}$, verifies the prefix
proof, log-inclusion proof, and signed tree head, and only on a verified owner slot issues a second
lookup for the shared warning slot. When both checks pass, the method returns \texttt{null}; it returns a
\texttt{MitmWarning} if the owner slot is absent or mismatched or the warning slot is occupied, and a
proof or head-anchoring failure throws. The live background monitor implements the
corresponding checks independently in \texttt{monitorConversation} (\cref{alg:detect}).

\subsection{Electron, Determinism, and Fail-Closed Behaviour}\label{ssec:apiref-eng-robust}

Three Electron obstacles arose. The bundler splits a chunk at every dynamic
import, so every dynamic import for the key-transparency code was removed and the whole layer now
lives in one bundle, with the networked path made opt-in and its failure caught separately. With
context isolation enabled, property assignments to \texttt{window} in the preload world are not visible
to the renderer, so objects injected through the console had no effect, whereas DOM custom events on
\texttt{window} cross the boundary. The application's typed event bus and those DOM custom events were
mismatched, resolved by dispatching on both. Several robustness issues appeared only end to end: a
server that generated a fresh signing key on each launch invalidated cached public keys, fixed by
deriving the server's keys deterministically from a seed; a panic on a freshly emptied log poisoned the
shared mutex, resolved by bootstrapping a genesis epoch at initialization; and an administrative reset
that rebuilds the server from its seed makes the demonstration repeatable.

On the client, the trusted computing
base is the verifier implementation, its cryptographic dependencies, and the directly implemented
ECVRF (\cref{ssec:impl-core-crypto}), together with the operator's two public keys, the STH signing
key and the VRF key. Because the prototype obtains these keys from the service during initialization,
their initial authenticity remains a trust assumption, and independent authenticated provisioning is a
production requirement. Subject to that assumption, the server, network, and log are not trusted, and
every response is re-checked. The verifier is fail-closed for integrity checks: any cryptographic or
proof-validation failure yields a reject verdict rather than falling through to acceptance. The present
key-fetch hook reports explicit mismatches and proof failures as a red warning but does not abort bundle
processing, so complete error-to-UI routing and session-blocking enforcement remain prototype
limitations (\cref{ch:discussion}). Production request handlers must also recover from poisoned locks
instead of relying only on the genesis-epoch bootstrap.
